%
% 3_Resultante.tex -- Abschnitt 3 Resultante 
%
% (c) 2026 Jakob Gierer, OST Ostschweizer Fachhochschule
%
% !TEX root = ../../buch.tex
% !TEX encoding = UTF-8
%
\section{Resultante
\label{resultante:section:resultante}}
\kopfrechts{Resultante}

Die Resultante ist ein Werkzeug der CAS, um das Vorhandensein gemeinsamer
Nullstellen zweier Polynome zu prüfen.
Die Resultante kann ebenfalls dazu verwendet werden, die Variablen eines
Gleichungssystems mit mehreren Variablen zu eliminieren.
Zu diesem Zweck wurde die Resultante im 19.~Jahrhundert untersucht
\cite{resultante:resultante}.
% In modernen CAS werden Resultanten unter anderem benutzt, um aus einer vorher bestimmten Gröbner-Basis auf die Lösung eines Gleichungssystems zu schliessen. 

\subsubsection{Nützliche Eigenschaften
\label{resultante:subsubsection:eigenschaften}}

Die Resultante bietet einige nützliche Eigenschaften.
Die erste und relevanteste Eigenschaft für dieses Kapitel ist, dass
die Resultante gemeinsame Nullstellen findet.
Zusätzlich dazu bietet sie auch einen Weg zur parametrischen Berechnung
der Nullstellen.

Wie bereits in (\ref{resultante:equation:detsyl}) gezeigt, ist die
Resultante bis auf ein Vorzeichen symmetrisch:
\begin{equation*}
    \Res(f,g)=(-1)^{mn}\Res(g,f).
\end{equation*}
Insbesondere wenn \(mn\) gerade ist, sind die beiden Resultanten gleich.
\index{Symmetrie}%

Eine weitere wichtige Eigenschaft ist die Elimination von Variablen,
was in der Anwendung in (\ref{resultante:equation:resz}) ebenfalls
genutzt wird.
\begin{beispiel}
    Die Elimination von Variablen wird am Beispiel der Polynome
    \(f(x,y)=x^2+xy+1\) und \(g(x,y)=x-y\) gezeigt.
    Hier treten nun zwei Unbekannte auf: \(x\) und \(y\).
    Für die Aufstellung der Resultante bedeutet dies, dass es zwei
    Möglichkeiten gibt, die Polynome und die Resultante zu betrachten:
    \begin{enumerate}
        \item Als Polynom in \(x\) mit Koeffizienten, die \(y\) enthalten
            \begin{equation*}
\def\l#1{\bgroup\color{gray!20}#1\egroup}
                \begin{aligned}
                    \Rightarrow\quad&\begin{minipage}[t]{8cm}\raggedright
                                $x$ wird zum ``Einsortieren'' der Koeffizienten
                                in der Sylvester-Matrix verwendet und somit
                                eliminiert.
                                \end{minipage}\\
                    \Rightarrow\quad& \Res_x(f,g)=
                    \det
                    \begin{pmatrix}
                        1 & \phantom{-}y & \phantom{-}1\\
                        1 & -y & \phantom{-}\l{0}\\
                        \l{0} & \phantom{-}1 & -y
                    \end{pmatrix}
                    =2y^2+1.
                \end{aligned}
            \end{equation*}
        \item Als Polynom in \(y\) mit Koeffizienten, die \(x\) enthalten
            \begin{equation*}
                \begin{aligned}
                \Rightarrow\quad& \Res_y(f,g)=
                \det
                \begin{pmatrix}
                    \phantom{-}x & x^2+1\\
                    -1 & x
                \end{pmatrix}
                =2x^2+1.
		\end{aligned}
\qedhere
            \end{equation*}
    \end{enumerate}
    % \kreis{1} Als Polynom in \(x\) mit Koeffizienten, die \(y\) enthalten
    % \begin{equation*}
    %     \begin{aligned}
    %         &\Longrightarrow x \phantom{x} \text{wird zum ``einsortieren'' der Koeffizienten in der }\\
    %         &\hspace{21pt}\text{Sylvester-Matrix verwendet und somit eliminiert.}\\
    %         &\Longrightarrow \Res_x(f,g)=\det(\syl_x(f,g))=2y^2+1
    %     \end{aligned}
    % \end{equation*}
    % \kreis{2} Als Polynom in \(y\) mit Koeffizienten, die \(x\) enthalten
    % \begin{equation*}
    %     \Longrightarrow \Res_y(f,g)=\det(\syl_y(f,g))=2x^2+1.
    % \end{equation*}
    \label{resultante:beispiel:elimination}
\end{beispiel}
\noindent
Das tiefgestellte \(x\) oder \(y\) gibt dabei an, nach welcher Resultante aufgelöst wird.
Dabei verschwindet diese Variable aus dem Resultanten-Polynom (\(R(y)=\Res_x(f,g)=2y^2+1\)).

Die Resultante liegt im von \(f\) und \(g\) erzeugten Ideal, das man als \((f,g)\) schreibt, also
\begin{equation}
    \Res(f,g) \in (f,g).
\end{equation}
Das bedeutet, es existieren Polynome \(s(X)\) und \(t(X)\), daraus folgt:
\begin{equation}
    \Res(f,g) = s(X)f(X)+t(X)g(X).
    \label{resultante:equation:kombi}
\end{equation}
Im Prinzip ähnelt dies dem Span aus der linearen Algebra.
Bei Vektoren ist
\begin{equation}
    \Span(v_1,v_2)=av_1+bv_2,
\end{equation}
wobei \(a\), \(b\) Zahlen sind.
Beim Ideal \((f,g)\) dürfen die Faktoren \(s\), \(t\) selbst Polynome sein.

\begin{beispiel}
    Am Beispiel von \(f(X)=X-2\) und \(g(X)=X-5\) wird gezeigt:
    Die Resultante aus diesen Polynomen ist
    \begin{equation*}
        \Res(f,g)=\det
        \begin{pmatrix}
            1 & -2\\
            1 & -5
        \end{pmatrix}
        =-3.
    \end{equation*}
    Dabei ist
    \begin{equation*}
        -1\cdot f+ 1\cdot g=-3,
    \end{equation*}
    mit den Polynomen \(s(X)=-1\) bzw. \(t(X)=1\) und deshalb
    \begin{equation*}
        \Res(f,g)\in(f,g). \qedhere
        \label{resultante:equation:ideal}
    \end{equation*}
    \label{resultante:beispiel:ideal}
\end{beispiel}

Eine ähnliche Form zu (\ref{resultante:equation:kombi}) erhält man
durch den erweiterten euklidischen Algorithmus.
Aus dieser Form lässt sich ein effizientes Berechnungsverfahren für
die Resultante ableiten, das Subresultanten-Verfahren
(siehe Abschnitt~\ref{resultante:section:fortsetzung}).
\index{Subresultanten-Verfahren}%
